% Ad Atticum 7.13 --- headnote
% ad-atticum-07-13, 24 January 49 BC, Minturnae

\textit{Cicero to Atticus, written from Minturnae on
the ninth day before the Kalends of February,
24 January 49 BC (the manuscript dateline:
\textit{Scr. Menturnis ix K. Febr. a. 705 (49)}). One
day after the Formian letter; the southward retreat
through Latium has reached the Liris. Labienus's defection
from Caesar has now arrived as confirmed news and is the
opening note of the letter.}

\textit{Section 1 is the political assessment hardening
into a verdict on the war. Labienus is a \emph{hērōs}
--- a piece of civil action more brilliant than any in
memory, if only because it gave Caesar pain. But the
\emph{genus belli} (kind of war) is what Cicero stops to
fix in his teeth: civil only in so far as it has been
born of one ruined citizen's audacity, not of any quarrel
between citizens. Caesar holds the City unprotected and
full of supplies; he treats temples and roofs as plunder,
not as fatherland; he cannot even pretend at anything
\emph{politikōs} --- in a statesmanlike way. The judgement
on Pompey is just as sharp: \emph{astratēgētos}, no
strategist --- a leader to whom Picenum was unfamiliar
country.}

\textit{Section 2 turns to the present hopelessness: two
treacherously-retained legions, an unwilling levy, the
time for terms lost; the ship at sea without a rudder.
Section 3 is the family triage --- whether to send the
two boys off to Greece, what to do with Tullia and
Terentia when one moment he fears barbarians coming to
the City and the next remembers Dolabella is there. The
question for Atticus is formally split: \emph{pros to}
(what is right in itself) and then with an eye to opinion,
lest his women be reproached for staying when good men
are leaving. Section 4 closes on Atticus's reading of the
news --- the closing scholarly joke about the ``riddle of
the Oppii from Velia'' being more obscure than the number
of Plato is Cicero's wave at the famously cryptic
\emph{nuptial number} of \emph{Republic} 8.}
